\documentclass[11pt,a4paper]{article}
\usepackage[margin=22mm]{geometry}
\usepackage[T1]{fontenc}
\usepackage{lmodern,amsmath,booktabs,tabularx,xcolor,hyperref,fancyhdr}
\hypersetup{colorlinks=true,urlcolor=blue,linkcolor=blue}
\setlength{\parindent}{0pt}\setlength{\parskip}{6pt}
\pagestyle{fancy}\fancyhf{}\fancyhead[L]{\small COE619 | Zahid Alabadllah | 202045160}\fancyfoot[C]{\thepage}
\newcommand{\hot}[1]{\par\textbf{Hot take:} #1\par}
\begin{document}
\section*{COE619 Project Understanding Guide}
Energy-Aware Predictive VM Consolidation in Cloud Data Centers\\
Advisor: Dr. Mohamed Faten Zhani | Evidence checked: 15 November 2026

\textbf{Purpose.} Explain the project through PR2 without memorizing long derivations. Each formula below keeps only what is needed to explain the implementation. The research question is whether simple prediction improves the balance between energy and performance. A positive answer has not yet been demonstrated.

\section{The project in everyday words}
A physical host is a server. A virtual machine (VM) is a software-defined computer sharing that server. Consolidation means putting VMs on fewer hosts so empty hosts can turn off. Live migration moves a running VM between hosts. It has a cost, so moving earlier is not automatically better.

\hot{An active but idle host still consumes 175 W in our model. Turning an empty host off saves more than merely lowering its CPU use.}

A cloudlet is a simulated task, not a real application deployment. A processing element (PE) represents a CPU core. MI means million instructions, an amount of work. MIPS means million instructions per second, a processing rate, not electrical power. RAM is working memory; bandwidth is network transfer capacity; storage is disk capacity. SLA means service level agreement. Our SLA scores are simulation proxies, not measured application response times.

\href{https://www.usenix.org/conference/nsdi-05/live-migration-virtual-machines}{Clark et al. (2005): live migration}. \href{https://javadoc.io/doc/org.cloudsimplus/cloudsim-plus/8.0.0/index.html}{CloudSim Plus: simulation entities}.

\section{What the code checks}
\textbf{Overload detection: CPU only.} Host utilization is allocated MIPS divided by physical host MIPS. Our host has 8 PEs at 1,000 MIPS each, so its capacity is 8,000 MIPS. Allocating 6,400 MIPS gives 80\% utilization.

\textbf{Placement: multiple resources.} A destination must have suitable CPU, RAM, bandwidth and storage. Detection asks whether a host should shed load; placement asks where a VM can fit. Forecast checks can also influence destination acceptance. A trigger is a policy decision, not proof that a migration completed.

\textbf{Requested CPU} is what a VM asks for. \textbf{Allocated CPU} is what it receives. A gap is unmet demand. The overload threshold (80\% for THR) is an early warning. Saturation accounting uses at least 99.9\% allocated host capacity. These are different quantities.

\hot{An 80\% warning gives headroom; workload jumps, migration time and sampling mean it cannot guarantee prevention of saturation.}

\newpage
\section{The three policies}
\textbf{THR: static threshold.} Flag overload when current utilization exceeds 0.80. It reacts to the present reading, not necessarily to a saturated server.

\textbf{IQR: interquartile range.} Sort the historical utilization readings. Q1 is the 25th percentile and Q3 is the 75th percentile. IQR is the spread of their middle half:
\[\mathrm{IQR}=Q3-Q1,\qquad \text{threshold}=1-0.75\,\mathrm{IQR}.\]
All quantities are fractions, not percentages. Example: IQR = 0.10 gives threshold = 0.925, or 92.5\%. More spread lowers the threshold. Less spread raises it. The implementation uses all active history entries, with static fallback if history is insufficient. The factor 0.75 is a chosen safety setting, not a proven optimum.
\hot{IQR changes the warning level according to past variability. It does not forecast the next reading.}
\paragraph{Reference.} Research convention: \href{https://doi.org/10.1002/cpe.1867}{Beloglazov and Buyya (2012)}; \href{https://beloglazov.info/papers/2012-optimal-algorithms-ccpe.pdf}{free full-paper PDF}. Implementation: \href{https://javadoc.io/doc/org.cloudsimplus/cloudsim-plus/8.0.0/index.html}{CloudSim Plus 8.0.0 API}. Research convention, not a formal standard.\par

\textbf{Our LR: ordinary least squares (OLS) linear regression.} It fits the straight line that minimizes squared errors over the last 10 active host-history readings. The fitted line has a slope (change in utilization per second) and an intercept (line height at time zero). We evaluate it at current time plus 300 seconds:
\[\text{forecast}=\text{fitted line evaluated 5 minutes ahead}.\]
\hot{Extend the recent trend forward. If current CPU or the forecast exceeds 80\%, flag overload.}
Before 10 readings exist it behaves like THR. Ten regularly spaced 5-minute readings span 45 minutes from first to last, and represent roughly 50 minutes of observations. Host-history events can be irregular, so the code uses actual timestamps; the window is not guaranteed to cover exactly 50 minutes. K=10 and the 5-minute horizon are initial choices, not tuned optima.

This is unweighted OLS, not Beloglazov's locally weighted regression (LOESS). Do not describe prediction itself as new to VM consolidation. Our contribution must be established through a careful implementation and evaluation. A lightweight formula also does not prove negligible execution overhead.

\textbf{Which VM moves?} All three policies use minimum migration time (MMT) selection. In this model, less RAM generally means less transfer time. Equal-size VMs limit how informative that choice is. Bandwidth and the simulator's migration settings affect transfer duration; this is not a packet-level implementation of real pre-copy migration.

\newpage
\section{Energy and performance scores}
\textbf{Power and energy.} For an active host:
\[P(u)=175+75u\quad\text{watts}.\]
Here u is CPU fraction: 0 means idle and 1 means fully used. Off-host power is modelled as zero. Watts measure current power; kWh measures energy accumulated over time. One kWh is 3,600,000 joules.
\hot{Idle is 175 W, full CPU is 250 W; running at 200 W for two hours consumes 0.4 kWh.}
These are chosen simplified model parameters, not measurements of our laptop or a calibrated hardware claim. Energy is integrated over recorded active-host history until simulation end, not only cloudlet execution. The configured horizon is 86,400 s; recorded end times can exceed it slightly. Cooling, network equipment and measured predictor energy are not included.

\textbf{SLATAH: SLA time per active host.}
SLATAH measures how much of a server's active time was spent at the saturation cutoff (at least 99.9\% CPU).
\[\text{One host's saturation fraction}=\frac{T_s}{T_a}.\]
\textbf{Ts} = seconds spent saturated. \textbf{Ta} = total seconds the host was active (powered on).

\textbf{What does ``ratio'' mean?} One amount divided by another. Here it tells us the share of active time spent saturated. Example: $60/600=0.10=10\%$. The host was saturated for 60 of its 600 active seconds.

\textbf{Our reported SLATAH:} calculate this fraction for each host with active time, then take their average. Each included host has equal weight, even if hosts ran for different lengths of time.
\hot{SLATAH answers: ``How much of the servers' active time was spent saturated?'' Lower is better.}
The code counts recorded intervals starting at the cutoff. A zero means no counted saturation; brief events may be missed, and zero does not prove migrations prevented overload.

\textbf{PDM: performance degradation due to migration.}
PDM measures the estimated share of processing work lost because VMs were moved.
\[\text{One VM's lost-work fraction}=\frac{0.10\times\text{VM capacity}\times\text{migration seconds}}{\text{work baseline}}.\]
The top estimates lost work in MI: 10\% of VM capacity (total MIPS across its cores), multiplied by migration seconds. The bottom is finished-task work in MI; if none finished, the code uses capacity times execution time instead.

\textbf{Example only:} 2 MI lost out of 100 MI gives $2/100=0.02=2\%$. Reported PDM averages the eligible VMs' fractions, including zero for VMs that did not migrate.
\hot{PDM answers: ``What share of work was estimated lost during migration?'' Lower is better.}
The 10\% is our model assumption, not the final PDM score. The fallback work baseline is an approximation, not measured total CPU demand.
\paragraph{Reference.} Research convention: \href{https://doi.org/10.1002/cpe.1867}{Beloglazov and Buyya (2012)}; \href{https://beloglazov.info/papers/2012-optimal-algorithms-ccpe.pdf}{free full-paper PDF}. Implementation: \href{https://javadoc.io/doc/org.cloudsimplus/cloudsim-plus/8.0.0/index.html}{CloudSim Plus 8.0.0 API}. Research convention, not a formal standard.\par

\newpage
\textbf{SLAV: composite SLA violation score.}
SLAV combines the saturation score (SLATAH) and the migration-loss score (PDM):
\[\mathrm{SLAV}=\mathrm{SLATAH}\times\mathrm{PDM}.\]
\textbf{Example only:} if SLATAH is 0.10 and PDM is 0.02, then SLAV is $0.10\times0.02=0.002$. This is a combined research score, not a measured percentage of customer requests that failed.
\hot{SLAV answers: ``What is the combined saturation-and-migration score?'' Lower is better, but a zero can hide a cost.}
In our results, SLATAH is zero, so SLAV is zero even when PDM is positive. Zero SLAV therefore does not mean there was no migration cost or performance loss.

\textbf{ESV: energy and SLA violations.} ESV multiplies energy in kWh by SLAV. It is energy scaled by the combined score. Since our SLAV is zero, ESV is also zero and cannot distinguish these runs.

\textbf{How to read these together:} SLATAH = saturation time; PDM = estimated work lost while moving; SLAV = their product. Also show energy, migrations, active-host hours and unmet demand. These are research conventions, not formal SLA compliance certificates.

\textbf{Unmet demand.} At roughly 300-second intervals, compare running VMs' requested and allocated MIPS. Sum positive shortfalls times 300 seconds and divide by sampled requested MIPS times 300 seconds.
\hot{What share of sampled requested CPU work was not supplied?}
The result is dimensionless. Multiplying by 100 expresses it as a percentage. This is coarse sampling, not exact continuous integration or application latency. A zero estimate can miss short events. The current code measures it; an older version hardcoded zero.

\textbf{Migration count} counts completed moves. \textbf{Active-host hours} sums all hosts' active seconds and divides by 3,600. Two active hosts for three hours contribute six host-hours. \textbf{Predictive trigger count} counts overload-check calls returning true from the forecast, including possible destination checks. Five triggers do not mean five distinct migrations.

\newpage
\section{Simple experimental design}
We compare THR, IQR and our OLS predictor on five fixed cases. The exact trace paths and VM order are saved in src/main/resources/workloads. Every policy receives the same list for a case. There is no runtime trace sampling, seed loop or averaging across cases.
\begin{center}\small\begin{tabular}{lrrl}\toprule Case & Hosts & VMs & Purpose\\\midrule
diurnal & 4 & 8 & Controlled day/night example\\
real a, real b, real c & 10 & 20 each & Three fixed real-trace selections\\
heavy & 6 & 16 & High-demand real-trace example\\\bottomrule\end{tabular}\end{center}
The real groups are three disjoint selections frozen once from the 1,052-trace pool, using a shuffled list with generator initialization 619 during manifest preparation. This is provenance, not a setting used when running the project. They are not claimed to be statistically representative. The heavy case uses the same 16 curated traces in a fixed name-sorted order. The full pool mean is 12.31\%; the heavy pool mean is 78.34\%, with minimum trace mean 68.12\%.

The experiment now answers: ``What happens on these clearly specified workloads?'' It does not establish an average improvement across all data centers. A fixed input list also cannot guarantee bit-identical simulation output: tie/iteration ordering can vary. Additional repetitions or workload days should be added only where needed to support a specific final claim.

\section{Current results, one row per case and policy}
The table comes from the current 15-row results/experiments.csv. There are no five-seed averages. Full SLA, unmet-demand and forecast-check fields are retained in the CSV and results/summary.md.
\begin{center}\small\begin{tabular}{llrrr}\toprule Case & Policy & kWh & Moves & Host-hours\\\midrule
diurnal & THR & 5.2488 & 2 & 25.67\\
diurnal & IQR & 5.2285 & 2 & 25.67\\
diurnal & LR & 5.7720 & 3 & 28.50\\
real a & THR & 21.9444 & 0 & 120.42\\
real a & IQR & 21.9444 & 0 & 120.42\\
real a & LR & 21.9444 & 0 & 120.42\\
real b & THR & 21.8661 & 0 & 120.42\\
real b & IQR & 21.8661 & 0 & 120.42\\
real b & LR & 21.8661 & 0 & 120.42\\
real c & THR & 21.5899 & 0 & 120.42\\
real c & IQR & 21.5899 & 0 & 120.42\\
real c & LR & 21.5899 & 0 & 120.42\\
heavy & THR & 8.2520 & 4 & 38.34\\
heavy & IQR & 7.0405 & 0 & 30.83\\
heavy & LR & 9.5788 & 5 & 45.39\\
\bottomrule\end{tabular}\end{center}
In this batch LR used more energy than THR in both diurnal and heavy cases. The three real groups produced no migrations and equal policy energy within each group. Heavy IQR used less energy and fewer active-host hours. These are case-specific observations, not universal rankings.\par Read the rows directly: lower kWh is less simulated energy; moves counts completed migrations; host-hours sums server active time. The predictor must earn its benefit against the baselines. A forecast flag is an overload-check result, not proof of a completed proactive migration. These results do not establish general predictive superiority.

Older seed-based runs are retained in results/archive/seed-workflow for provenance. Their numbers belong to an older input-selection/order design and must not be mixed with this table.
\newpage
\section{What to run and what the folders mean}
Traces are input text files, not programs. Each value is CPU percentage for a five-minute sample. diurnal holds synthetic day/night traces; real/20110303 holds all 1,052 real traces; heavy holds the 16 high-demand traces. These three folders support current experiments. The unused active subset has been removed. The top-level 20110303 folder contains four historical samples, unused by the current runner.

The workloads folder contains only lists of paths, not duplicate trace data. real a, b and c name the fixed groups in those lists. Seeds and averages are no longer part of the current workflow. They remain in the historical experiment archive, where a seed controlled random selection/order and an average combined repeated runs.

\textbf{From the repository's VS Code PowerShell terminal:}
\begin{verbatim}
./scripts/run.ps1 -Scenario diurnal
./scripts/run.ps1 -Scenario heavy
./scripts/run.ps1 -Scenario real_a
./scripts/run.ps1
\end{verbatim}
The last command runs all five cases and all three policies, giving 15 result rows. real b and real c can also be run by name. Demo output goes to results/demo. Use -OutputDir results only to deliberately replace the report evidence and then update the tables. The helper builds in target/meeting-build to keep command-line output separate from the IDE's compiled files.

\textbf{Code roles.} ExperimentRunner reads the lists and compares policies. ConsolidationExperiment builds the simulation and calculates metrics. PredictiveLinearRegressionPolicy fits the OLS line. Baseline programs preserve PR1 history.

\textbf{Advisor answers.} ``I compare three policies on the same fixed workloads. Energy measures the cost of running hosts; migrations measure movement; SLA and unmet-demand fields describe modelled performance impact. Prediction has not established a general advantage. Next I will validate the current approach before expanding it.''
\end{document}
